% Kandydat restrukturyzacji; nie jest aktywnym rozdziałem Atlasu.
\section{Transformata dystrybucji o zwartym nośniku}
\providecommand{\atlasref}[2]{#2}

\begin{twierdzenie}[Transformata dystrybucji o zwartym nośniku]
\label{thm:compact-fourier22}
Jeżeli $T\in\mathcal E'(\R^d)$ ma zwarty nośnik $K$,
to jego transformata Fouriera jest dystrybucją regularną
reprezentowaną przez gładką funkcję
\[
F_T(\xi)=(2\pi)^{-d/2}\langle T_x,e^{-ix\cdot\xi}\rangle.
\]
Dla pewnego $m\in\N_0$ i każdego wielowskaźnika $\beta$
istnieje $C_\beta>0$ takie, że
$|\partial_\xi^\beta F_T(\xi)|
\le C_\beta(1+|\xi|)^m$
dla wszystkich $\xi\in\R^d$.
W działaniu na funkcji wykładniczej używamy
przedłużenia z Lematu~\ref{lem:compact-action21}.
\end{twierdzenie}
\begin{proof}
Wybierzmy $\chi\in\mathcal D$ równą $1$ na
otoczeniu $K$. Dla $\beta\in\N_0^d$ funkcje
$x\mapsto\chi(x)(-ix)^\beta e^{-ix\cdot\xi}$
mają nośniki w ustalonym zwartym zbiorze $\supp\chi$.
Na zwartym otoczeniu ustalonego parametru $\xi$
wszystkie pochodne względem $x$ i parametru są ciągłe
i jednostajnie ciągłe na iloczynie tych dwóch zbiorów.
Stosując twierdzenie Newtona--Leibniza do ilorazów
różnicowych względem parametru w seminormach
$\mathcal D_{\supp\chi}$, dostajemy
\[
 \partial_\xi^\beta F_T(\xi)
 =(2\pi)^{-d/2}
 \langle T_x,(-ix)^\beta e^{-ix\cdot\xi}\rangle.
\]
Ciągłość tych pochodnych wynika z tego samego argumentu.
Z Lematu~\ref{lem:compact-action21} wybieramy $C,m,L$
niezależne od $\beta,\xi$, gdzie $L$ jest zwarty i
$|T(f)|\le C\max_{|\alpha|\le m}\sup_{x\in L}
|\partial^\alpha f(x)|$.
Dla $|\alpha|\le m$ wzór Leibniza daje
\[
 \partial_x^\alpha\bigl((-ix)^\beta e^{-ix\cdot\xi}\bigr)
 =\sum_{\gamma\le\alpha}\binom\alpha\gamma
 \partial_x^\gamma((-ix)^\beta)
 (-i\xi)^{\alpha-\gamma}e^{-ix\cdot\xi}.
\]
Ponieważ $x\in L$, pochodne wielomianu są ograniczone,
a moduł czynnika wykładniczego dla rzeczywistych $\xi$
wynosi $1$. Ponadto $|\xi^{\alpha-\gamma}|
\le(1+|\xi|)^m$. Po zsumowaniu skończonej liczby
składników otrzymujemy $C_\beta(1+|\xi|)^m$ jako
oszacowanie modułu $\partial_\xi^\beta F_T(\xi)$.
W szczególności $T_{F_T}$ jest temperowana.

Aby sprawdzić równość dystrybucji, ustalmy
$\varphi\in\mathcal S$ i $c_d=(2\pi)^{-d/2}$.
Dla $R\ge1$ połóżmy
\[
 J_R(x)=c_d\chi(x)\int_{[-R,R]^d}
 e^{-ix\cdot\xi}\varphi(\xi)\,\mathrm d\lambda_d(\xi).
\]
Sumy Riemanna zbiegają do $J_R$ w
$\mathcal D_{\supp\chi}$: dla każdego rzędu pochodnych
względem $x$ wynika to z jednostajnej ciągłości funkcji
podcałkowej i jej pochodnych na iloczynie zwartych
zbiorów $\supp\chi$ i $[-R,R]^d$.
Stąd
\[
 \langle T,J_R\rangle=
 \int_{[-R,R]^d}F_T(\xi)\varphi(\xi)
 \,\mathrm d\lambda_d(\xi).
\]
Dla każdego $q\in\N_0$ wzór Leibniza daje
\[
 p_q(J_R-\chi\mathcal F\varphi)
 \le C_{\chi,q,d}\int_{\R^d\setminus[-R,R]^d}
 (1+|\xi|)^q|\varphi(\xi)|\,\mathrm d\lambda_d(\xi)
 \longrightarrow0.
\]
Pochodne pod całką istnieją na mocy całkowalności
każdej wagi wielomianowej razy $|\varphi|$.
Wszystkie funkcje po lewej mają nośniki w $\supp\chi$,
więc zbieżność zachodzi w jednym etapie $\mathcal D$.
Z ciągłości $T$ i bezwzględnej całkowalności
$F_T\varphi$ dostajemy
\[
 \langle\mathcal FT,\varphi\rangle
 =\langle T,\chi\mathcal F\varphi\rangle
 =\int_{\R^d}F_T(\xi)\varphi(\xi)
 \,\mathrm d\lambda_d(\xi).
\]
Pierwsza równość wykorzystuje jednoznaczne przedłużenie
$T$ na $\mathcal S$ i regułę odcięcia.
Zatem $\mathcal FT=T_{F_T}$.
\end{proof}
